%% Class-M local inputs, formal recurrences, and residue comparisons.

\providecommand{\Vir}{\mathrm{Vir}}
\providecommand{\cW}{\mathcal W}
\providecommand{\ClaimStatusProvedHere}{}
\providecommand{\ClaimStatusConditional}{}
\providecommand{\ClaimStatusDefinitional}{}
\providecommand{\ClaimStatusOpen}{}

\chapter{Class-M Residue Towers: Local Data and Comparison Problems}
\label{chap:shadow-tower-other-class-m-platonic}

\section*{Three mathematical layers}

For every algebra in this chapter, the local OPE belongs to
Beilinson level~\(1\).  A scalar residue coefficient belongs to a
completed ordered bar model at level~\(2\).  An asymptotic law belongs
to a uniform family over parameter space.  We write
\[
 \mathsf H_{\mathrm{form}},\qquad
 \mathsf H_{\mathrm{res}},\qquad
 \mathsf H_{\mathrm{asymp}}
\]
for the corresponding recurrence, residue, and uniform-asymptotic
packages.  Their separation permits exact local calculations to serve
as inputs while the geometric comparisons remain visible proof
obligations.

\section{A typed formal master equation}
\label{sec:class-m-master-equation}

\begin{proposition}[Formal quadratic recurrence]
\label{prop:shadow-tower-master-equation-classM}
\ClaimStatusProvedHere
Fix a line \(\ell\) in a chosen strong-generator presentation and
formal initial data \(R_2,R_3,R_4\) in a coefficient field \(F\).
The package \(\mathsf H_{\mathrm{form}}(\cA,\ell)\) defines
\begin{equation}
 R_r(\cA,\ell)=
 -\frac{1}{rK_\ell}
 \sum_{\substack{j+k=r+2\\3\le j\le k<r}}
 \epsilon_{j,k}\,jkR_j(\cA,\ell)R_k(\cA,\ell),
 \qquad
 \epsilon_{j,k}=\begin{cases}1,&j<k,\\[2pt]\frac12,&j=k,
 \end{cases}
 \label{eq:master-eqn-classM}
\end{equation}
whenever \(K_\ell\in F^\times\).  It determines a unique formal
sequence in \(F\).
\end{proposition}

\begin{proof}
At arity \(r\), every index in the sum is smaller than \(r\).
Induction therefore gives existence and uniqueness.
\end{proof}

\begin{remark}[Relation with the Virasoro recurrence]
\label{rem:vir-master-eqn-proof-relation}
The Virasoro recurrence is one instance of
Proposition~\ref{prop:shadow-tower-master-equation-classM}.  A residue
theorem for \(\cA\) consists of a contraction that identifies the
transferred Maurer--Cartan coefficients with this formal sequence.
\end{remark}

\begin{remark}[Ordered provenance]
\label{rem:e1-ordered-provenance}
The residue package is built in the ordered \(E_1\) bar construction.
Symmetric averaging is a later chain map.  The comparison records the
ordering, collision strata, signs, and scalar projection.
\end{remark}

\begin{remark}[Linewise universality]
\label{rem:line-universality-classM}
A common recurrence across two generator lines follows from an
intertwining map between their transferred scalar \(L_\infty\) models.
The local equality of selected OPE coefficients supplies the first
input to that map.
\end{remark}

\section{The \(W_3\) algebra: two local lines}
\label{sec:w3-two-lines}

In the standard normalization, \(W_3\) is strongly generated by the
Virasoro field \(T\) and a primary field \(W\) of weight~\(3\).  The
exact local identities include
\[
 T(z)W(w)\sim\frac{3W(w)}{(z-w)^2}
 +\frac{\partial W(w)}{z-w},
\]
and the \(WW\)-OPE contains the quasi-primary
\(\Lambda={:}TT{:}-\frac{3}{10}\partial^2T\) with coefficient
\(32/(22+5c)\).  The Virasoro Gram determinant therefore supplies the
same divisor \(c(5c+22)\) to the local \(T\)- and \(W\)-channel
calculations.

\subsection{The \(T\)-line}

\begin{proposition}[Virasoro subalgebra comparison]
\label{prop:w3-tline-virasoro-inheritance}
\ClaimStatusConditional
\emph{Type signature:}
\textup{(Open quadrant, completed ordered residue presentation,
level~\(2\), compatible \(W_3\) and Virasoro residue packages).}
If the inclusion generated by \(T\) intertwines the ordered
differentials, contractions, and scalar projections, then
\[
 S_r(W_3,T;\mathsf H_{\mathrm{res}}^{W_3})
 =S_r(\Vir_c;\mathsf H_{\mathrm{res}}^{\Vir})
\]
for every arity defined by the two packages.
\end{proposition}

\begin{proof}
The \(T\)-generated vertex subalgebra is \(\Vir_c\).  The stated
intertwining carries its transferred Maurer--Cartan element and scalar
contraction through the inclusion.
\end{proof}

\begin{corollary}[Formal \(T\)-line asymptotic]
\label{cor:w3-tline-asymptotic}
\ClaimStatusConditional
If the compared formal sequence satisfies the Virasoro weighted
recurrence and \(\mathsf H_{\mathrm{asymp}}\) identifies it uniformly
with the residue sequence, then
\[
 \lim_{c\to\infty}c^{r-2}S_r(W_3,T)
 =\frac{8(-6)^{r-4}}r
\]
on the domain of that asymptotic package.
\end{corollary}

\subsection{The \(W\)-line}

\begin{proposition}[Formal \(W\)-line sequence]
\label{prop:w3-wline-closed-form}
\ClaimStatusConditional
Fix the standard \(WW\)-OPE normalization and a scalar projection of
the \(W\)-word subcomplex.  Under
\(\mathsf H_{\mathrm{form}}(W_3,W)\), Equation~\eqref{eq:master-eqn-classM}
defines the full formal \(W\)-line sequence.  Under the compatible
\(\mathsf H_{\mathrm{res}}\), this sequence is the geometric
\(W\)-line residue tower.
\end{proposition}

\begin{proposition}[Leading \(W\)-line scale]
\label{prop:w3-wline-leading-asymptotic}
\ClaimStatusOpen
The leading large-\(c\) scale of the geometric \(W\)-line is the
uniform limit obtained from the explicit \(WW\)-channel contraction.
Its proof obligation is the construction of
\(\mathsf H_{\mathrm{asymp}}(W_3,W)\) and control of the divisor
\(5c+22\).
\end{proposition}

\begin{remark}[Two scales]
\label{rem:w3-two-scales}
The \(T\)-line begins with the central term of the Virasoro OPE\@.  The
\(W\)-line begins with the sixth-order \(WW\) pole and the
\(\Lambda\)-channel.  Their asymptotic normalizations are therefore
separate inputs.
\end{remark}

\begin{remark}[Denominator geometry]
\label{rem:w3-wline-denominator-pattern}
Powers of \(5c+22\) arise from inversion of the level-four Virasoro
Gram form.  A residue denominator theorem further records the number
of inverse-Gram insertions in each collision tree.
\end{remark}

\begin{remark}[Reference normalization]
\label{rem:w3-wline-ref-to-exponential-base}
An exponential base for the \(W\)-line is meaningful after the
generator normalization, scalar contraction, and large-\(c\) scaling
have been fixed.
\end{remark}

\section{Bershadsky--Polyakov reduction}
\label{sec:bp-shadow-tower}

The universal Bershadsky--Polyakov algebra has generators
\(J,G^+,G^-,T\) and standard central charge
\[
 c_{\mathrm{BP}}(k)=-\frac{(2k+3)(3k+1)}{k+3}.
\]
The shifted conformal vector used in a second convention has scalar
\[
 c_{\mathrm{BP}}^{\mathrm{shift}}(k)
 =2-\frac{24(k+1)^2}{k+3}
 =-\frac{2(12k^2+23k+9)}{k+3}.
\]
These are distinct local presentations.

\begin{proposition}[Formal BP \(T\)-line]
\label{prop:bp-tline-rational}
\ClaimStatusConditional
After choosing one conformal vector, substitution of its central
charge into the formal Virasoro recurrence gives a rational sequence
in \(k\).  A compatible residue contraction identifies this sequence
with the BP \(T\)-line tower.
\end{proposition}

\begin{proposition}[Heisenberg \(J\)-line]
\label{prop:bp-jline-gaussian}
\ClaimStatusConditional
For a normalized Heisenberg current \(J\), Wick's theorem computes the
local vacuum correlators by pairings.  A residue tower concentrated on
this Gaussian line is obtained when the chosen contraction factors
through those pairings.
\end{proposition}

\begin{corollary}[Central-charge reflection datum]
\label{cor:bp-tline-koszul-conductor}
\ClaimStatusProvedHere
The two rational functions displayed above determine their reflection
sums by direct substitution once a level involution \(k\mapsto k^{!}\)
has been fixed.  The standard and shifted conformal vectors yield
separate conductor conventions.
\end{corollary}

\begin{remark}[Graded \(G^+G^-\) channel]
\label{rem:bp-gpm-graded-riccati}
The half-integral generator line carries parity and charge.  Its formal
recurrence is therefore a graded transferred equation rather than a
scalar substitution in the \(T\)-line recurrence.
\end{remark}

\section{A universal asymptotic comparison problem}
\label{sec:universal-asymptotic-factor}

\begin{theorem}[Linewise asymptotic transport]
\label{thm:universal-asymptotic-factor}
\ClaimStatusConditional
Let \(X_r(K_\ell)=K_\ell^{r-2}S_r(\cA,\ell)\).  Suppose a family line admits
$\mathsf H_{\mathrm{form}}+\mathsf H_{\mathrm{res}}$,
\(X_4(K_\ell)\to2\), and, for every \(r\ge5\),
\[
 X_r(K_\ell)=-\frac{6(r-1)}rX_{r-1}(K_\ell)+E_r(K_\ell),
 \qquad E_r(K_\ell)\longrightarrow0
\]
locally uniformly on the parameter domain.  Then
\[
 \lim_{K_\ell\to\infty}K_\ell^{r-2}S_r(\cA,\ell)
 =\frac{8(-6)^{r-4}}r.
\]
\end{theorem}

\begin{proof}
The assertion holds at \(r=4\).  Passing to the limit in the displayed
recurrence and applying induction gives
\[
 A_r=-\frac{6(r-1)}rA_{r-1},\qquad A_4=2.
\]
Its unique solution is \(A_r=8(-6)^{r-4}/r\).
\end{proof}

\begin{corollary}[Family specializations]
\label{cor:universal-asymptotic-factor}
\ClaimStatusConditional
Every family whose line satisfies the hypotheses of
Theorem~\ref{thm:universal-asymptotic-factor} has the same normalized
leading sequence.
\end{corollary}

\begin{remark}[Scope of universality]
\label{rem:class-m-asymptotic-universality}
Universality is a theorem about compared transferred scalar models.
The generator weight, parity, additional channels, and family
singularities enter through the hypothesis package.
\end{remark}

\section{Principal \(W_N\) families}
\label{sec:wn-shadow-tower}

\begin{proposition}[Generator-line decomposition]
\label{prop:wn-line-decomposition}
\ClaimStatusProvedHere
The standard strong generating set of the principal \(W_N\)-algebra
has conformal weights \(2,3,\ldots,N\).  Its ordered bar complex is
therefore multigraded by generator words, and its transferred scalar
model decomposes into line and mixed-channel components after a
compatible projection has been chosen.
\end{proposition}

\begin{proposition}[\(T\)- and higher-spin lines]
\label{prop:wn-tline-and-w-lines}
\ClaimStatusConditional
The \(T\)-line compares with \(\Vir_c\) under the same subcomplex
hypothesis as Proposition~\ref{prop:w3-tline-virasoro-inheritance}.
Each higher-spin line requires its own OPE Gram form and residue
contraction.
\end{proposition}

\begin{remark}[Mixed channels]
\label{rem:wn-mixed-line-channels}
Collision trees containing several strong-generator weights form the
mixed-channel sector.  Their contribution is computed by the full
multigraded transfer.
\end{remark}

\section{The \(W_\infty[\mu]\) family}
\label{sec:w-infinity-shadow}

\begin{proposition}[Filtered line decomposition]
\label{prop:w-infinity-line-decomposition}
\ClaimStatusConditional
At a parameter value where a completed strong-generator presentation
and continuous residue contraction exist, the tower decomposes into a
completed sum of spin lines and mixed channels.
\end{proposition}

\begin{remark}[The Virasoro line]
\label{rem:w-infinity-tline-recovers-virasoro}
The spin-two subalgebra supplies the local Virasoro line.  Equality of
geometric residue coefficients follows from the compatible completed
subcomplex comparison.
\end{remark}

\begin{theorem}[Large-parameter asymptotic target]
\label{thm:shadow-tower-w-infinity-asymptotic}
\ClaimStatusOpen
The \(W_\infty[\mu]\) asymptotic problem asks for a parameter domain,
a completed residue package, and uniform estimates in spin and arity.
These data define the large-parameter limit.
\end{theorem}

\begin{remark}[Parameter dependence]
\label{rem:w-infinity-mu-dependence-open}
The structure constants and truncation ideals vary with \(\mu\).
Uniformity is consequently formulated on a specified domain avoiding
the relevant discriminant locus.
\end{remark}

\begin{remark}[Closed-form problem]
\label{rem:c-w-infinity-closed-form-open}
A closed form is the solution of the completed multispin transferred
equation with its convergence estimates.
\end{remark}

\section{Super-Yangian lines}
\label{sec:super-yangian-shadow}

\begin{proposition}[Even/odd line decomposition]
\label{prop:super-yangian-fermionic-line}
\ClaimStatusProvedHere
A homogeneous generating set of a super vertex algebra splits the
ordered bar model by parity.  Koszul signs decorate every collision
tree containing odd generators.
\end{proposition}

\begin{proposition}[Formal super-Yangian recurrence]
\label{prop:super-yangian-tline-shadow}
\ClaimStatusConditional
After fixing the completed presentation, parity conventions, and
scalar contraction, the transferred equation defines even, odd, and
mixed formal sequences.  A residue comparison identifies them with
geometric coefficients.
\end{proposition}

\begin{proposition}[Conductor coordinate]
\label{prop:super-yangian-kappa}
\ClaimStatusConditional
The scalar conductor of a super-Yangian line is the reflection sum of
the specified central or supertrace coordinate.  Its value follows
from that coordinate and the chosen duality involution.
\end{proposition}

\begin{corollary}[Asymptotic target]
\label{cor:super-yangian-tline-asymptotic}
\ClaimStatusOpen
The normalized large-parameter limit is determined by the graded
recurrence together with uniform control of odd and mixed channels.
\end{corollary}

\begin{remark}[Fermionic sign package]
\label{rem:super-yangian-fermionic-sign-open}
The sign package consists of an ordering convention, a suspension
convention, and the induced signs on the transferred brackets.
\end{remark}

\section{Triplet \(W(p)\) Cartan line}
\label{sec:wp-triplet-cartan-shadow}

\begin{proposition}[Cartan-line computation through arity six]
\label{prop:wp-cartan-shadow-through-r6}
\ClaimStatusOpen
For each \(p\), the arities \(2\le r\le6\) are obtained from the
explicit triplet OPE, its vacuum Gram matrices, and an ordered residue
contraction.  The open proof obligation is the family-uniform
construction of these data and their scalar comparison.
\end{proposition}

\begin{remark}[Small parameters]
\label{rem:wp-cartan-values-p2-p3}
The cases \(p=2,3\) provide the first worked computations once the
normalizations of the Cartan generator and invariant pairing are
fixed.
\end{remark}

\begin{remark}[Quartic cross-channel]
\label{rem:wp-cross-channel-quartic}
At arity four, the pure Cartan line and the first mixed collision
channel are separated by the multigraded transfer.
\end{remark}

\begin{remark}[Scope]
\label{rem:wp-scope-honesty}
The proposition records a finite computation programme.  An all-arity
triplet tower is a further theorem about the completed logarithmic
category and its residue model.
\end{remark}

\section{Status ledger}
\label{sec:class-m-closed-form-summary}

\begin{table}[htbp]
\centering
\begin{tabular}{c|c|c|c}
family & exact local input & residue comparison & asymptotic package\\
\hline
\(W_3\) & \(TT,TW,WW\) OPE and Gram form & conditional & open\\
BP & chosen conformal vector and OPE & conditional & open\\
\(W_N\) & strong-generator presentation & conditional & open\\
\(W_\infty[\mu]\) & parameterwise OPE & open & open\\
super-Yangian & graded presentation & conditional & open\\
\(W(p)\) & finite local OPE data & open & open
\end{tabular}
\caption{Epistemic status of the class-M family lanes.}
\label{tab:class-m-shadow-summary}
\end{table}

\begin{table}[htbp]
\centering
\begin{tabular}{c|c|c}
lane & formal object & geometric input\\
\hline
spin two & Virasoro recurrence & compatible residue subcomplex\\
higher spin & line recurrence & line Gram form and contraction\\
mixed & multigraded transfer & full collision complex
\end{tabular}
\caption{Formal and geometric components of the completed spin decomposition.}
\label{tab:w-infinity-family-integer}
\end{table}

\section{Frontier problems}
\label{sec:frontier-questions-shadow-tower-other-class-m}

The class-M programme has three concrete obligations: construct the
ordered residue package for each presentation, prove its comparison
with the formal transferred equation, and establish uniform estimates
on the chosen parameter domain.  These steps turn the exact local OPE
data into geometric residue towers and then into asymptotic theorems.
